% ====================================================================
\section{Error analysis}\label{app:errors}
% ====================================================================

\paragraph{Fits.}
I fitted all lines by least squares with equal weights (\texttt{curve\_fit} from SciPy). The exception is the joint
fit of the two \ce{MnCl2} scans in \cref{sec:mncl2}, where I weighted each scan, and each part of the search scan,
by the inverse of its own residual scatter. I scaled the covariance matrix with the residual variance
$s^2=\sum_ir_i^2/(N-n_\mathrm p)$ ($N$ data points, $n_\mathrm p$ fit parameters), so that the uncertainties reflect
the actual scatter. For averages over several scans I took the larger of the internal and external uncertainty:
\begin{equation}\label{eq:wmean}
  \bar x=\frac{\sum_jw_jx_j}{\sum_jw_j},\qquad w_j=\sigma_j^{-2},\qquad
  \sigma_{\bar x}=\max\!\left[\Bigl(\sum_jw_j\Bigr)^{-1/2},\ \frac{s_x}{\sqrt n}\right],
\end{equation}
where $s_x$ is the standard deviation of the $n$ individual values.

\paragraph{Reference field.}
By \cref{eq:resonance}, $B_\text{DPPH}=h\nu/(g_\text{DPPH}\muB)$, so
\begin{equation}
  \frac{\sigma_B}{B}=\sqrt{\left(\frac{\sigma_g}{g}\right)^2+\left(\frac{\sigma_\nu}{\nu}\right)^2}
  =\sqrt{(\num{1.0e-4})^2+(\num{2.3e-6})^2}=\num{1.0e-4},
\end{equation}
and $\sigma_B=\qty{0.033}{mT}$. Only the literature value of $g_\text{DPPH}$ counts; the frequency contributes nothing.

\paragraph{Field scale.}
With $B(U)=k\,B_\text{Hall}(U)$ and $k=B_\text{DPPH}/B_\text{Hall}(U_\text{D})$ (\cref{eq:fieldscale}),
\begin{equation}
  \frac{\sigma_k}{k}=\sqrt{\left(\frac{\sigma_{B_\text{DPPH}}}{B_\text{DPPH}}\right)^2
  +\left(\frac{B_\text{Hall}'(U_\text{D})\,\sigma_{U_\text{D}}}{B_\text{Hall}(U_\text{D})}\right)^2}
  =\num{1.4e-4},
\end{equation}
where $\sigma_{U_\text{D}}=\qty{1.9}{mV}/\sqrt{11}$. I converted widths with the slope
$B'=k\,B_\text{Hall}'(U_\text{D})=\qty{60.76}{mT/V}$. Different descriptions of the calibration curve
(polynomial degree, scale or zero-point error of the display) give a slope uncertainty of
$\sigma_{B'}/B'=\qty{3.0}{\percent}$. The field lag (\cref{sec:field}) is not included; it enters as a one-sided
systematic.

\paragraph{Line width and \texorpdfstring{$T_2$}{T2}.}
With $\gamma=g\muB/\hbar$ of the respective sample, $\Bpp=2B'G/\sqrt3$ and $T_2=1/(\gamma B'G)$. Both carry the
relative uncertainties of $G$ and $B'$:
\begin{equation}
  \frac{\sigma_{T_2}}{T_2}=\frac{\sigma_{\Bpp}}{\Bpp}=\sqrt{\left(\frac{\sigma_G}{G}\right)^2+\left(\frac{\sigma_{B'}}{B'}\right)^2}.
\end{equation}
Here $\sigma_G$ includes the fit error and the model uncertainty from \cref{sec:dpph}.

\paragraph{\texorpdfstring{$g$}{g}-factor relative to DPPH.}
For a line at $U_x$ and frequency $\nu_x$,
\begin{equation}\label{eq:grel}
  g_x=\frac{h\nu_x}{\muB\,k\,B_\text{Hall}(U_x)}
  =g_\text{DPPH}\,\frac{\nu_x}{\nu_\text{DPPH}}\,\frac{B_\text{Hall}(U_\text{D})}{B_\text{Hall}(U_x)} .
\end{equation}
The factor $k$ drops out, and with it the scale error of the probe. The remaining uncertainties are those of the
literature value $g_\text{DPPH}$, the reference position $U_\text{D}$ and the line position $U_x$:
\begin{equation}
  \frac{\sigma_{g_x}}{g_x}=\sqrt{\left(\frac{\sigma_{g_\text{DPPH}}}{g_\text{DPPH}}\right)^2
  +\left(\frac{B'_\text{Hall}(U_\text{D})\,\sigma_{U_\text{D}}}{B_\text{Hall}(U_\text{D})}\right)^2
  +\left(\frac{B'_\text{Hall}(U_x)\,\sigma_{U_x}}{B_\text{Hall}(U_x)}\right)^2}.
\end{equation}
For \ce{MnCl2}, $\sigma_{U_\text{D}}=\qty{12}{mV}$ (\cref{eq:udmn}) contributes \num{0.0046} and the fit error of $U_x$
contributes \num{0.0024}. I estimated the other contributions in \cref{sec:mncl2} (line shape \num{0.0076}, sample in
the cavity \num{0.0052}, shape of the calibration curve and $g_\text{DPPH}$ together \num{0.0010}) from how much $g$
changes when I vary the line model, the calibration curve and the corrections. All contributions added in quadrature
give \num{0.0106}.

\paragraph{Guide wavelength.}
The phase rotation $\Delta\varphi$ for the plunger displacement $\Delta z$ gives $\lamg=\ang{720}\,\Delta z/|\Delta\varphi|$
and $\nu_\mathrm c=\nu\sqrt{1-(\lambda_0/\lamg)^2}$ with $\lambda_0=c/\nu$. Error propagation gives
\begin{equation}
  \frac{\sigma_{\lamg}}{\lamg}=\sqrt{\left(\frac{\sigma_{\Delta\varphi}}{\Delta\varphi}\right)^2
  +\left(\frac{\sigma_{\Delta z}}{\Delta z}\right)^2},\qquad
  \sigma_{\nu_\mathrm c}=\frac{c^2}{\nu_\mathrm c\,\lamg^3}\,\sigma_{\lamg} .
\end{equation}
Compared with the phase, the plunger scale (\qty{0.01}{mm}) contributes nothing. Because $\lamg$ is inversely
proportional to $|\Delta\varphi|$, the symmetric interval $\Delta\varphi=\qty{-88(13)}{\degree}$ becomes the asymmetric
interval $\lamg=46^{+8}_{-6}\,\unit{mm}$. I computed its limits directly from $|\Delta\varphi|=\ang{75}$ and \ang{101}.

% ====================================================================
\section{Fit results of the DPPH scans}\label{app:fits}
% ====================================================================
% generated by analysis/tables.py -- do not edit by hand
\begin{center}
  \footnotesize\setlength{\tabcolsep}{3.4pt}\captionsetup{hypcap=false}
  \captionof{table}{Fit results of all analyzed DPPH scans. Convolution model (\cref{eq:convolution}): position $U_0$, intrinsic
  HWHM $G$, length $\lambda$ of the field distribution, mixing angle $\varphi_\ast=\operatorname{atan2}(d,a)$. Simple
  Lorentzian fit: $G_\mathrm{L}$, $\varphi_\mathrm{L}$. Last column: ratio of the residual rms of both fits.
  \qty{1}{mV} corresponds to \qty{60.76}{\micro T}; attenuators in dB; arrow: sweep direction.}
  \label{tab:dpphfits}
  \sisetup{table-number-alignment=center, uncertainty-mode=compact}
  \begin{tabular}{@{}l S[table-format=1.2] S[table-format=2.0] S[table-format=2.0] c S[table-format=3.1] S[table-format=1.4] S[table-format=1.2(2)] S[table-format=2.1(2)] S[table-format=-3.0] S[table-format=1.2] S[table-format=-3.0] S[table-format=1.1]@{}}
    \toprule
    Scan & {Mod.\ (V)} & {X382A} & {X375A} & & {$S_\mathrm{pp}$ (\si{\micro V})} & {$U_0$ (V)} & {$G$ (mV)} & {$\lambda$ (mV)}
      & {$\varphi_\ast$ (\si{\degree})} & {$G_\mathrm{L}$ (mV)} & {$\varphi_\mathrm{L}$ (\si{\degree})} & {$\frac{\mathrm{rms_L}}{\mathrm{rms}_\ast}$} \\
    \midrule
    07b & 3 & 1 & 2 & $\uparrow$ & 190.1 & 4.9617 & 3.55(4) & 15.8(5) & 141 & 7.64 & -172 & 3.3 \\
    08 & 2 & 1 & 2 & $\uparrow$ & 189.4 & 4.9635 & 3.77(5) & 14.1(5) & 143 & 7.76 & -173 & 3.2 \\
    09 & 1 & 1 & 2 & $\uparrow$ & 105.3 & 4.9448 & 3.64(8) & 14.5(10) & 143 & 7.53 & -170 & 3.0 \\
    10 & 0.5 & 1 & 2 & $\uparrow$ & 52.8 & 4.9487 & 3.50(8) & 14.5(9) & 144 & 7.41 & -169 & 2.7 \\
    11 & 0.25 & 1 & 2 & $\uparrow$ & 26.8 & 4.9495 & 3.64(7) & 14.4(7) & 146 & 7.56 & -168 & 3.3 \\
    12 & 0.1 & 1 & 2 & $\uparrow$ & 10.4 & 4.9501 & 3.52(4) & 15.2(7) & 145 & 7.43 & -167 & 3.4 \\
    13 & 1 & 0 & 2 & $\uparrow$ & 115.9 & 4.9457 & 3.94(9) & 12.3(11) & 156 & 7.51 & -160 & 2.3 \\
    14 & 1 & 3 & 2 & $\uparrow$ & 65.3 & 4.9461 & 3.64(9) & 12.3(10) & 161 & 7.03 & -153 & 2.3 \\
    15 & 1 & 6 & 2 & $\uparrow$ & 28.9 & 4.9462 & 3.42(7) & 10.8(7) & 138 & 6.65 & -179 & 3.1 \\
    16 & 1 & 10 & 2 & $\uparrow$ & 39.0 & 4.9467 & 3.49(5) & 12.1(5) & 85 & 7.12 & 131 & 3.9 \\
    16b & 1 & 10 & 2 & $\uparrow$ & 27.9 & 4.9477 & 3.49(5) & 11.4(6) & 84 & 7.00 & 128 & 3.4 \\
    17b & 1 & 20 & 2 & $\uparrow$ & 30.4 & 4.9484 & 3.55(7) & 12.2(7) & 54 & 7.12 & 101 & 3.0 \\
    18 & 1 & 6 & 2 & $\downarrow$ & 41.5 & 4.9188 & 3.17(5) & 10.9(6) & 124 & 6.01 & 180 & 2.4 \\
    24 & 1 & 0 & 6 & $\uparrow$ & 50.4 & 4.9486 & 3.84(9) & 12.8(11) & 145 & 7.37 & -170 & 2.6 \\
    25 & 1 & 0 & 10 & $\uparrow$ & 22.0 & 4.9494 & 3.77(7) & 12.9(10) & 143 & 7.30 & -171 & 2.9 \\
    26 & 1 & 0 & 15 & $\uparrow$ & 7.4 & 4.9491 & 3.75(8) & 12.5(9) & 145 & 7.21 & -170 & 2.9 \\
    27 & 1 & 0 & 20 & $\uparrow$ & 2.1 & 4.9498 & 3.67(5) & 13.8(8) & 143 & 7.25 & -168 & 3.6 \\
    20 & 1 & 0 & 2 & $\uparrow$ & 93.9 & 4.9498 & 4.02(11) & 13.0(9) & 68 & 8.02 & 113 & 2.8 \\
    \bottomrule
  \end{tabular}
\end{center}

